\section{Step 30: Conjoined Intrinsic Discriminator}

\paragraph{Deflationary start.}
Step 29 narrowed the neutral P2 carrier to 1,851 atomic rewrite packages while keeping the reference support. Step 30 conjoins a second intrinsic predicate, strict closure-currency minimization.

\paragraph{Currency predicate.}
For each Step-29 survivor, the C1/C4 currency is computed from the support's internal audit-vector cancellation work:
\[
(\#\text{active audit axes},\ \ell_1\text{ shadow price},\ \#\text{nonzero audit terms},\ \text{charge span}).
\]
The discriminator is the lexicographic minimum of this currency tuple. This is not a multiplet-count criterion and does not use the removed exterior/GUT prior.

\paragraph{Computation.}
The Step-29 survivor count is reproduced as 1,851. The strict currency minimum is \(2|4|4|2\), leaving 4 conjoined survivors. The reference support has currency \(6|546|20|10\), so it fails the strict currency minimum.

\paragraph{Verdict.}
The verdict is \texttt{TYPED\_NO\_GO}. Strict closure-currency minimization is neutral and intrinsic, but it selects tiny low-price packages rather than the reference support.

\paragraph{Next grammar delta.}
The next intrinsic predicate should be F24/F47 role-obstruction or P3 protocol-holonomy on the Step-29 atomic packages.

\paragraph{Scope.}
This is a finite-carrier diagnostic. It is not a physical gauge-structure theorem, not a real-world selection mechanism, and not a frame-transfer status upgrade.
